\section{Introducci\'on y contexto del dominio}
\label{sec:introduccion}

En la provincia de Los R\'ios y su zona de influencia agr\'icola en Quevedo, las asociaciones de productores de cacao administran centros de acopio encargados de recibir grano en baba y grano seco (variedades CCN-51 y Nacional Arriba), evaluar su calidad f\'isica, consolidar lotes para comercializaci\'on y liquidar los pagos a los socios. En la operaci\'on cotidiana, una parte considerable de estos registros se lleva en cuadernos f\'isicos o planillas aisladas, lo que origina errores en el destare de sacos, discrepancias en los descuentos por humedad o impurezas, retrasos en el cruce de anticipos de insumos y falta de trazabilidad al emitir liquidaciones.

El presente documento desarrolla la actividad formativa \texttt{FORM-SIM-01} (registrada en el SGA como \texttt{GA-SIM-01 -- FORM}) como ejercicio preparatorio estructurado para \texttt{TA-IND-02}. Sobre el dominio asignado por sorteo (N.\textsuperscript{o} 6: \textit{Asociaci\'on de productores de cacao: acopio y pagos}), se formulan 20 requisitos (12 funcionales y 8 de calidad cuantificados) conforme a ISO/IEC/IEEE 29148:2018 \cite{iso29148} e ISO/IEC 25010:2023 \cite{iso25010}, y se aplican tres t\'ecnicas complementarias de priorizaci\'on: MoSCoW, Modelo de Kano y votaci\'on ponderada de 100 puntos \cite{pohl2025,wiegers2013,ireb2026}.

El objetivo t\'ecnico es contrastar la priorizaci\'on efectuada por el analista frente a las preferencias de las partes interesadas del centro de acopio, documentando coincidencias, discrepancias y una propuesta sustentada para la primera versi\'on del sistema.
